% =====================================================================
% Hito 2 - Avance (máximo 2 páginas, sin portada, con nombres)
% EYP2307 Análisis de Regresión, 2026-2
%
% Cómo usar en Overleaf: subir esta carpeta completa (main.tex + figuras/).
% Cada sección indica qué rol la escribe y cuánto espacio aproximado tiene.
% Las figuras se guardan en figuras/ con el nombre indicado; mientras no
% existan, aparece un recuadro en su lugar (el documento compila igual).
% Las cifras de la Tabla 1 salen de outputs/tables/flujo_muestral.csv.
% =====================================================================
\documentclass[10pt]{article}
\usepackage[utf8]{inputenc}
\usepackage[T1]{fontenc}
\usepackage[spanish,es-tabla,es-nodecimaldot]{babel}
\usepackage[a4paper,margin=1.5cm]{geometry}
\usepackage{lmodern}
\usepackage{graphicx}
\usepackage{booktabs}
\usepackage{float}   % [H] = la figura o tabla queda donde se escribe
\usepackage[font=small,labelfont=bf]{caption}
\usepackage[hidelinks]{hyperref}
\usepackage{enumitem}
\setlist{nosep}
\setlength{\parskip}{2pt}
\setlength{\parindent}{0pt}
% Deja cortar líneas largas (p. ej. la URL) sin salirse del margen.
\setlength{\emergencystretch}{3em}
% Títulos de sección más compactos (para caber en 2 páginas).
\makeatletter
\renewcommand{\section}{\@startsection{section}{1}{0pt}%
  {-1.2ex plus -0.5ex}{0.5ex}{\normalfont\large\bfseries}}
\makeatother
% Menos espacio alrededor de figuras y tablas.
\setlength{\textfloatsep}{6pt plus 2pt}
\setlength{\intextsep}{6pt plus 2pt}
\setlength{\abovecaptionskip}{3pt}

% Figura con recuadro de reemplazo mientras el archivo no exista.
\newcommand{\figura}[2]{%
  \IfFileExists{#1}{\includegraphics[width=#2]{#1}}%
  {\fbox{\parbox[c][4cm][c]{0.95#2}{\centering\small Falta \texttt{#1}}}}}

\begin{document}

\begin{center}
  {\large\bfseries Modelo de regresión para la pérdida dado el incumplimiento (LGD)\\
  en créditos hipotecarios de tasa fija originados en 2007}\\[2pt]
  EYP2307 Análisis de Regresión --- Avance (Hito 2)\\
  % Verificar la grafía de los nombres (Entrega 1 dice "Amelie"; el boceto, "Amelia").
  \textbf{Integrantes:} Diego Olivares, Amelie Ramírez, Enzo Balboni y Joaquín Ortuzar
\end{center}

% ---------------------------------------------------------------------
\section{Introducción y datos}
% Rol 4. BORRADOR dejado por el Rol 2 (Diego), redactado a partir de la propuesta
% del Hito 1 (Entrega 1 y boceto). El Rol 4 lo revisa, escribe la motivación y lo termina.

Este trabajo modela la pérdida dado el incumplimiento (LGD) en créditos hipotecarios de
tasa fija originados en 2007 y adquiridos por Freddie Mac. La pregunta es: \emph{¿qué
características de originación del crédito y de la propiedad explican la LGD de estas
hipotecas que terminaron en pérdida hasta marzo de 2026?}
\textit{[Rol 4: una o dos líneas de motivación, escritas por el grupo.]}

\textbf{Datos.} Se usa la muestra de créditos originados en 2007 del \emph{Single-Family
Loan-Level Dataset} de Freddie Mac, de acceso público y gratuito (con registro) en
\url{https://www.freddiemac.com/research/datasets/sf-loanlevel-dataset}.
La muestra tiene
dos archivos: originación (50.000 créditos, uno por fila) y desempeño mensual (cerca de
3 millones de registros), sin encabezado, delimitados por ``|'' y unidos por el
identificador del crédito (\emph{Loan Sequence Number}). Las variables y sus códigos se
definen en la \emph{General User Guide} de Freddie Mac (julio de 2026).

\textbf{Respuesta y predictores.} La respuesta es
$\mathrm{LGD} = \textit{Actual Loss} / \textit{Zero Balance Removal UPB}$, es decir, la
fracción del saldo adeudado al liquidar el crédito que se pierde. La población son los
créditos con \emph{Zero Balance Code} 02, 03, 09 o 15 (venta a terceros, venta corta o
castigo, venta de la propiedad recuperada y venta de cartera) y con \emph{Actual Loss}
informado. Los predictores candidatos son variables conocidas al otorgar el crédito:
puntaje de crédito, LTV, CLTV, DTI, tasa de interés, porcentaje de seguro hipotecario y
monto original (continuos); y propósito del préstamo, ocupación, tipo de propiedad, número
de unidades, región, canal, primera vivienda y número de deudores (categóricos). No se usan
los componentes de la pérdida (precio de venta, gastos, recuperaciones del seguro), porque
forman parte del cálculo de la respuesta (fuga de información).

\textbf{Frecuencia y periodo.} La unidad de observación es el crédito (corte transversal),
cuyo desempeño se registra con frecuencia mensual. Los créditos se originaron en 2007 y su
desempeño se observa hasta marzo de 2026; las liquidaciones con pérdida ocurrieron entre
2008 y 2025. El eje temporal es el mes de liquidación (\emph{Zero Balance Effective Date}),
que se usa para la validación: el 20\,\% de los créditos liquidados más recientemente forma
la muestra de validación. Se excluyen las liquidaciones de los tres meses previos al corte,
para las que Freddie Mac aún no informa la pérdida.

% ---------------------------------------------------------------------
\section{Limpieza de datos}
% Rol 1 (párrafo) + Tabla 1. ~6 líneas de texto.

\textit{[Rol 1: párrafo de limpieza. Puntos: códigos ``no disponible'',
casos completos y advertencia sobre DTI, divisiones censales y agrupaciones, partición
temporal. Detalle en cleaning\_log.md.]}

\begin{table}[H]
  \centering\small
  \caption{Flujo muestral.}
  \begin{tabular}{lrr}
    \toprule
    Paso & Quedan & Excluidos \\
    \midrule
    Créditos en la muestra 2007                                  & 50.000 & --- \\
    Cierre con pérdida (Zero Balance Code 02, 03, 09, 15)        &  4.484 & 45.516 \\
    Actual Loss informado                                         &  4.043 &    441 \\
    Sin liquidaciones en los 3 meses previos al corte (202603)    &  4.043 &      0 \\
    Sin faltantes en los predictores candidatos                   &  3.960 &     83 \\
    Estados con división censal (sin Puerto Rico)                 &  3.956 &      4 \\
    \midrule
    \quad Entrenamiento (liquidación hasta 2014-12)               &  3.153 & --- \\
    \quad Validación (liquidación desde 2015-01)                  &    803 & --- \\
    \bottomrule
  \end{tabular}
\end{table}

% ---------------------------------------------------------------------
\section{Análisis exploratorio}
% Rol 2. Figura 1 (varios paneles en una sola imagen) + ~8 líneas.
% Guardar como figuras/figura1.pdf (o .png, cambiando la extensión abajo).

\begin{figure}[H]
  \centering
  \figura{figuras/figura1.pdf}{\linewidth}
  \caption{\textit{[Rol 2: descripción de los paneles.]}}
\end{figure}

\textit{[Rol 2: hallazgos del EDA, faltantes y atípicos, transformaciones evaluadas
para los predictores y comparación de la LGD de los créditos sin DTI.]}

% ---------------------------------------------------------------------
\section{Modelo preliminar y diagnóstico}
% Rol 3. Tabla 2 + Figura 2 (paneles de diagnóstico) + ~8 líneas.

\begin{table}[H]
  \centering\small
  \caption{Modelos OLS anidados ajustados con entrenamiento ($n = 3.153$); error de
  validación contra la LGD original ($n = 803$).}
  \begin{tabular}{lcccccc}
    \toprule
    Modelo & Parámetros & $R^2$ ajustado & AIC & BIC & RMSE val. & MAE val. \\
    \midrule
    Completo & --- & --- & --- & --- & --- & --- \\
    Reducido & --- & --- & --- & --- & --- & --- \\
    \midrule
    \multicolumn{7}{l}{Prueba F (reducido vs.\ completo): $F = $ ---, $p = $ ---} \\
    \bottomrule
  \end{tabular}
\end{table}

\begin{figure}[H]
  \centering
  \figura{figuras/figura2.pdf}{\linewidth}
  \caption{\textit{[Rol 3: residuos vs.\ ajustados, QQ-plot, residuos vs.\ fecha de
  liquidación, apalancamiento y distancia de Cook.]}}
\end{figure}

\textit{[Rol 3: comparación de modelos, diagnóstico (supuestos, VIF, atípicos e
influyentes) y resultado de la validación.]}

% ---------------------------------------------------------------------
\section{Conclusiones y próximos pasos}
% Rol 4 (con el grupo). ~4 líneas.

\textit{[Rol 4: conclusiones breves del avance y qué falta para el póster.]}

\end{document}
